La respuesta inicial de la IA resultó sintácticamente válida y modeló adecuadamente la carga inicial del carrito y los dos bloques opcionales (\texttt{opt}) para modificar la cantidad y eliminar un artículo. Pese a ello, se identificaron y corrigieron los siguientes aspectos:
\begin{itemize}[leftmargin=*,nosep]
\item \textbf{Inconsistencia de endpoints REST respecto al glosario:} La IA generó las rutas en inglés \texttt{GET /api/cart}, \texttt{PATCH /api/cart/items/\{itemId\}} y \texttt{DELETE /api/cart/items/\{itemId\}}. Se sustituyeron por las rutas del glosario común: \texttt{GET /api/carrito}, \texttt{PATCH /api/carrito/lineas/idLinea (cantidad)} y \texttt{DELETE /api/carrito/lineas/idLinea}.
\item \textbf{Alucinación de atributo inexistente (\texttt{color}):} En el primer bloque \texttt{opt}, la IA indicó \emph{``Cambia la cantidad de una prenda (talla/color)''}. En el modelo Entidad-Relación (\texttt{PRODUCTO}, \texttt{STOCK\_TALLA}, \texttt{LINEA\_CARRITO}) las variantes se gestionan exclusivamente por \texttt{talla}, por lo que se eliminó la referencia al color.
\item \textbf{Alineación con los métodos de las clases \texttt{Carrito}, \texttt{LineaCarrito} y \texttt{StockTalla}:} Se sustituyeron los mensajes en lenguaje natural del Backend por llamadas explícitas a los métodos de \texttt{clases.mmd}: \texttt{findCarritoConLineas(idCarrito)}, \texttt{Carrito.calcularSubtotal()}, \texttt{calcularEnvio()}, \texttt{calcularTotal()}, \texttt{StockTalla.hayStock(cantidad)}, \texttt{Carrito.modificarCantidad(idLinea, cantidad)}, \texttt{LineaCarrito.calcularImporte()} y \texttt{Carrito.eliminarLinea(idLinea)}, además de añadir la directiva \texttt{mirrorActors: false}.
\end{itemize}
